\documentclass[11pt]{article}

\usepackage{amsmath,amssymb,amstext,amsthm}
\usepackage{geometry}
\usepackage{fancyhdr}
\usepackage[pdftex]{graphicx}
\usepackage{xcolor}

\geometry{
  verbose,
  dvips,
  width=430pt, marginparsep=0pt, marginparwidth=0pt,
  top=90pt, headheight=12pt, headsep=20pt, footskip=30pt, bottom=80pt}

\setlength{\parskip}{\medskipamount}
\setlength{\parindent}{0pt}

\linespread{1.05}

\newtheorem{theorem}{Theorem}
\newtheorem{lemma}[theorem]{Lemma}
\newtheorem{proposition}[theorem]{Proposition}
\newtheorem{corollary}[theorem]{Corollary}
\newtheorem{conjecture}[theorem]{Conjecture}
\newtheorem{obs}{Observation}
\newtheorem{clm}{Claim}
\newtheorem{ques}{Question}
\newtheorem{problem}{Problem}

\newtheorem{ex}{Example}


\newcommand{\spc}{\hspace{0.08in}}
\newcommand{\vs}{\vspace*{.1in}}
\newcommand{\E}{\mathbb{E}}
\newcommand{\Prob}{\mathbb{P}}
\newcommand{\edit}[1]{\emph{\textcolor{blue}{#1}}}


\renewenvironment{proof}{{\noindent \textbf{\textit{Proof.}}}}
{\hfill $\Box$ \vspace*{0.1in}}
\newenvironment{proofc}{{\noindent \textit{Proof of Claim.}}}
{\hfill $\Box$}


\begin{document}

\begin{LARGE}\begin{center}\bf 13.2 Derivatives and Integrals of Vector Functions\end{center}
\end{LARGE}

\vs\vs



\section{Warm Up}

\textbf{Question 1}: What is the definition of the derivative?

\textbf{Problem 1:} Find the derivative of $\cos{x}$, $\ln{x}$, $xe^x$, $(1+\sqrt{x})^3$

\section{Motivation}

We have defined vector functions which seem useful to describe motion and curves/surfaces in 3-space. It seems natural to wonder if we have a notion of the derivative, which also helps us analyze the motion of functions?




\section{Definitions \& Theorems}

\begin{enumerate}
\item  If $r$ is a vector function, the derivative of $r$ denoted $r'$ is defined the same way we define it in calc I.

\begin{equation*}
    r'(t) = \lim_{h \to 0}\frac{r(t+h)-r(t)}{h}.
\end{equation*}



If this limit exists, geometrically it represents the \textbf{tangent vector} to the curve defined by $r$. It is a direction vector to the line tangent to the curve at a given $t$.

\item  This tangent vector will be useful to know and thus we define a \textbf{unit tangent vector} $T(t)$ as follows.

\begin{equation*}
    T(t) = \frac{r'(t)}{|r'(t)|}.
\end{equation*}

\item To compute the derivative for these vector functions we utilize the following theorem.

\begin{theorem}
    If $r(t) = \langle f(t),g(t),h(t) \rangle$ where $f,g,h$ are differentiable functions then $r'(t) = \langle f'(t), g'(t), h'(t) \rangle$. 
\end{theorem}

\item Many of the rules when taking the derivatives of functions also hold when taking the derivative of vector functions. These include the following.

\begin{itemize}
    \item $(u(t)+v(t))' = u'(t) + v'(t)$
    \item $(cu(t))' = cu'(t)$
    \item $(f(t)u(t))' = f'(t)u(t)+f(t)u'(t)$
    \item $(u(t)\cdot v(t))' = u(t)\cdot v'(t)+u'(t) \cdot v(t)$
    \item $(u(t) \times v(t))' = u'(t)\times v(t) + u(t) \times v'(t)$
    \item $(u(f(t))') = f'(t)u'(f(t))$
\end{itemize}

\item If we can find the derivative of a vector function which gives us a direction vector which is tangent to our space curve, we can use that, along with a point (Point of tangency) to find the \textbf{equation of the tangent line}! Just like we did in 12.5.

\item Similarly we can find the \textbf{integral} of a vector functions as follows.

\begin{equation*}
    \int_a^b r(t) = \left(\int_a^b f(t)\right) i + \left(\int_a^b g(t)\right) j + \left(\int_a^b h(t)\right) k.
\end{equation*}
   
   \end{enumerate}



\section{Examples}

\begin{problem}
  Find $r'(t)$ if $r(t) = \langle 1/t, 1/t^2, 1/t^3 \rangle$.
\end{problem}

\vspace{6cm}

\begin{problem}
   Find $r'(t)$ if $r(t) = e^t\sin(t)i - 2^tj + (1+\ln{t})^3k$.
\end{problem}

\vspace{7 cm}


\begin{problem}
Find $r''(t)$ in problem 1.
\end{problem}


\vspace{8cm}

\begin{problem}
Find $\int r(t) $ where $r(t) = \sqrt{t}i + 2t/(t^2+5)j - 4k$.
\end{problem}

\vspace{12cm}



\begin{problem}
    If $u(t) = \langle 3t^2, \sqrt[3]{t}, -5t \rangle$ and $v(t) = \langle 2t+3, -t^2, t-t^2 \rangle$, find $(u(t)\cdot v(t))'$ and $(u(t) \times v(t))'$.
\end{problem}
\newpage

\section{Scratch Work}

\end{document} 